El laboratorio estudia la conexión entre adquisición de datos en un ESP32
y distribución de mensajes mediante MQTT. La guía organiza experiencias
de conectividad, telemetría, control e integración serial
\cite{guia}.

\subsection{Marco teórico}
MQTT utiliza el modelo publicador--suscriptor: un cliente publica en un
tópico y el servidor distribuye mensajes según las suscripciones.
MQTT 3.1.1 requiere un transporte ordenado y bidireccional, como TCP.
Su nivel QoS~0 ofrece entrega como máximo una vez; no garantiza que cada
mensaje alcance al receptor \cite{mqtt311}. Por ello, observar una conexión
no basta para cuantificar pérdidas o recepción de lecturas.

En la API Arduino-ESP32, \texttt{analogRead} devuelve un código ADC crudo
sin calibrar, mientras \texttt{analogReadMilliVolts} proporciona un resultado
calibrado en milivoltios \cite{esp32adc}. Un código ADC obtenido de un
divisor con LDR no equivale directamente a iluminancia en lux. La relación
con iluminación requiere caracterizar sensor, circuito y condiciones de
referencia; aquí no se realiza esa conversión.

PubSubClient 2.8 selecciona MQTT 3.1.1 por defecto. Su sobrecarga
\texttt{publish(topic,payload)} construye un mensaje sin retención y sin
bits QoS~1/2; \texttt{loop()} atiende la comunicación del cliente
\cite{pubsub28}. Estas propiedades describen la biblioteca y permiten
interpretar el sketch de CP2; no sustituyen una captura del receptor.

\subsection{Antecedente de investigación y ubicación del trabajo}
Hunkeler, Truong y Stanford-Clark propusieron MQTT-S para extender la
comunicación publicador--suscriptor a redes de sensores con restricciones
de energía y ancho de banda \cite{mqtts}. El antecedente muestra la
adaptación del modelo a otros transportes. El presente laboratorio usa
MQTT sobre una conexión WiFi/TCP; no implementa MQTT-S ni reproduce sus
evaluaciones. Esta revisión breve aporta un antecedente pertinente y una
norma de referencia; no pretende una revisión exhaustiva del estado del arte.
